\documentclass[12pt]{article}

% -------------------------------------------------------
% PACKAGES
% -------------------------------------------------------
\usepackage[margin=1in]{geometry}
\usepackage{setspace}
\doublespacing
\usepackage{times}
\usepackage{amsmath,amssymb}
\usepackage{graphicx}
\usepackage{booktabs}
\usepackage{multirow}
\usepackage{array}
\usepackage{xcolor}
\usepackage{hyperref}
\hypersetup{colorlinks=true, linkcolor=blue, citecolor=blue, urlcolor=blue}
\usepackage{natbib}
\bibliographystyle{plain}
\usepackage{caption}
\usepackage{subcaption}
\usepackage{float}
\usepackage{microtype}
\usepackage{lineno}
\linenumbers % Adds line numbers for the reviewers

\begin{document}
\setcounter{page}{1} % GUARANTEES THE PAPER STARTS ON PAGE 1

% -------------------------------------------------------
% ANONYMISED TITLE 
% -------------------------------------------------------
\begin{center}
{\large \textbf{Dataset-Naive Cross-Validation Inflates ceRNA Network
Classifier Performance: Evidence from a Patient-Specific Tripartite
Graph Framework with Topology Sensitivity Diagnostics}}
\end{center}

\vspace{0.5cm}

% -------------------------------------------------------
% STRUCTURED ABSTRACT
% -------------------------------------------------------
\noindent\textbf{Abstract}

\noindent\textbf{Purpose:}
Multi-cohort liquid biopsy studies routinely report AUROC values of
0.90--0.95 for competing endogenous RNA (ceRNA) network classifiers, yet
few stratify cross-validation folds by dataset membership. We
investigate whether dataset-naive cross-validation systematically
inflates apparent performance by allowing the model to learn
platform-specific batch signals rather than cancer biology.

\noindent\textbf{Methods:}
We assembled a meta-cohort of 5,654 blood-based RNA samples from four
independent GEO datasets (GSE73002, GSE115513, GSE126094, GSE101684).
For each patient, a personalised directed tripartite
(circRNA--miRNA--mRNA) subgraph was constructed from experimentally
validated interactions (miRTarBase~v9.0, CircInteractome). Twenty-one
features were extracted per patient: four expression statistics and
seventeen graph-theoretic descriptors including Betti numbers and
persistence entropy from topological data analysis (TDA). An XGBoost
classifier was evaluated under two protocols: (i)~dataset-naive
stratified 5$\times$3-fold nested cross-validation (standard practice),
and (ii)~dataset-stratified nested cross-validation using
\texttt{StratifiedGroupKFold}, which holds out entire datasets in the
outer loop. A topology sensitivity analysis using Kruskal--Wallis tests
quantified the proportion of graph features showing significant
distributional shift across data sources.

\noindent\textbf{Results:}
Dataset-naive CV yielded AUROC~0.927 (95\%~CI: 0.921--0.934).
Dataset-stratified CV yielded AUROC~0.45--0.76 depending on the
held-out cohort, a reduction of up to 0.48~AUROC~points. Sensitivity
analysis revealed that 57.1\% (12/21) of topology features exhibit
significant cross-dataset distributional shift ($p < 0.001$,
Kruskal--Wallis), including community count, modularity, diameter,
spectral entropy, and all TDA features. Leave-one-dataset-out (LODO)
validation on the largest independent cohort (GSE73002, $n=4{,}113$)
yielded AUROC~0.749 (95\%~CI: 0.733--0.762), confirming that honest
cross-platform performance is substantially below within-cohort
estimates.

\noindent\textbf{Conclusions:}
Published ceRNA network classifiers that employ dataset-naive
cross-validation on multi-cohort data likely overestimate cancer
detection performance by learning platform identity rather than
regulatory biology. We provide a reproducible diagnostic pipeline
quantifying topology batch sensitivity and demonstrate that proper
dataset-stratified evaluation is necessary for credible performance
claims in multi-cohort liquid biopsy research.

\vspace{0.3cm}
\noindent\textbf{Keywords:} ceRNA network, liquid biopsy, cross-validation bias,
batch effects, topological data analysis, pan-cancer detection,
network topology, persistent homology

\newpage

% -------------------------------------------------------
% 1. INTRODUCTION
% -------------------------------------------------------
\section{Introduction}

Cancer remains a leading cause of global mortality, with 19.3~million
new cases and approximately 10~million deaths estimated in
2020~\citep{sung2021}. The survival disparity between early- and
late-stage diagnosis is stark: five-year survival for Stage~IV solid
tumours falls below 15\% for several major cancer types, whereas
Stage~I survival frequently exceeds 90\%. Liquid biopsy---the analysis
of tumour-derived circulating material recovered from peripheral
blood---offers a minimally invasive route to early
detection~\citep{cohen2018}.

Circular RNAs (circRNAs) have attracted sustained research attention as
liquid biopsy analytes. Covalently closed non-coding transcripts
generated by back-splicing events~\citep{memczak2013,hansen2013},
circRNAs are resistant to exonucleolytic degradation, accumulate stably
in blood exosomes and platelets, and exhibit tissue- and tumour-specific
expression patterns~\citep{kristensen2019,vo2019}. Their primary
functional role as competing endogenous RNAs (ceRNAs)---molecular
sponges that sequester microRNAs (miRNAs) and thereby de-repress miRNA
target messenger RNAs (mRNAs)~\citep{salmena2011}---positions them as
central regulatory nodes whose disruption is now recognised as a
hallmark of malignancy~\citep{meng2017}.

Computational classifiers built on ceRNA network features have reported
impressive AUROC values of 0.85--0.95 across multiple cancer
types~\citep{pan2026,yuan2026}. However, these studies share a common
methodological pattern: cross-validation folds are stratified by class
label but not by dataset membership. When multiple GEO cohorts are
combined into a single meta-cohort and folds are assigned randomly, the
model encounters samples from the same sequencing platform in both
training and test partitions. Because RNA expression profiles carry
strong platform-specific signatures---driven by differences in RNA
extraction, library preparation, sequencing depth, and circRNA
enrichment protocols---a model trained in this way may learn to
distinguish platforms as a proxy for distinguishing cancer from
healthy. The resulting AUROC reflects platform classification
performance as much as biological signal.

This confound has been documented in the broader multi-omics literature
\citep{leek2010} but has not been systematically quantified in the
ceRNA liquid biopsy subfield. The present work fills this gap by:
(i)~constructing a patient-specific tripartite ceRNA graph framework
and demonstrating the performance gap between dataset-naive and
dataset-stratified cross-validation; (ii)~providing a topology
sensitivity diagnostic that quantifies the proportion of network
features encoding dataset-specific batch effects; and
(iii)~establishing dataset-stratified cross-validation as the necessary
standard for credible multi-cohort ceRNA classifier evaluation.

% -------------------------------------------------------
% 1.1 Related Work
% -------------------------------------------------------
\subsection{Related Work}

Expression-based classifiers represent each patient as a continuous
vector of RNA abundance values. Pan~et~al.~\citep{pan2026} employed
nested cross-validation with Random Forest ensembles for small cell lung
cancer (SCLC) detection, achieving AUROC~0.950 in a single-institution
cohort. Yuan~et~al.~\citep{yuan2026} integrated three-dimensional
genome architecture with support vector machines for acute myeloid
leukaemia prediction. While these single-cohort studies are internally
consistent, multi-cohort studies face the dataset leakage problem
described above.

Network-based classifiers model biomolecular interactions explicitly.
Barab\'{a}si and Oltvai~\citep{barabasi2004} established that biological
networks exhibit characteristic topological properties---scale-free
degree distributions, modular community structure, hub-node
dominance---that are functionally meaningful and irreducible to
individual node properties. Population-average ceRNA networks have been
extensively constructed for prognostic biomarker discovery in uterine
corpus endometrial carcinoma~\citep{fan2025} and clear cell renal cell
carcinoma~\citep{farias2024}. Single-sample network (SSN) methodologies
such as LIONESS~\citep{kuijjer2019} estimate individual-specific
regulatory networks from aggregate transcriptomic data, but have
primarily been applied to prognostic survival tasks rather than
diagnostic liquid biopsy classification.

No prior study has (i)~evaluated a patient-specific ceRNA graph
classifier under dataset-stratified cross-validation, (ii)~quantified
the proportion of ceRNA topology features that encode dataset-specific
batch effects, or (iii)~reported the resulting performance gap as a
methodological warning for the field.

% -------------------------------------------------------
% 1.2 Contributions
% -------------------------------------------------------
\subsection{Contributions}

This paper makes three primary methodological contributions:

\begin{enumerate}
  \item A demonstration that dataset-naive cross-validation inflates
  AUROC by up to 0.48~points relative to dataset-stratified evaluation
  in a multi-cohort ceRNA liquid biopsy meta-cohort ($n=5{,}654$).

  \item A topology sensitivity diagnostic pipeline
  (Kruskal--Wallis cross-dataset feature distribution analysis) that
  identifies which graph-theoretic features encode platform batch
  effects versus biological signal, finding that 57.1\% of ceRNA
  topology features are significantly batch-sensitive ($p < 0.001$).

  \item A reproducible methodological benchmark establishing
  dataset-stratified \texttt{StratifiedGroupKFold} nested
  cross-validation as the appropriate evaluation standard for
  multi-cohort computational biomarker classifiers.
\end{enumerate}

% -------------------------------------------------------
% 2. MATERIALS AND METHODS
% -------------------------------------------------------
\section{Materials and Methods}

\subsection{Cohort Assembly}

Raw circRNA and miRNA expression count data were retrieved from the Gene
Expression Omnibus (GEO)~\citep{barrett2013} under four accession
numbers: GSE73002~\citep{li2018}, GSE115513, GSE126094, and GSE101684.
These four studies provide publicly accessible raw count data for
circRNA and miRNA expression in whole blood or blood-derived exosome
samples from both cancer patients and healthy volunteers, and
collectively span multiple cancer types and geographical populations.
After quality filtering---removal of samples with more than 20\%
missing expression values across profiled transcripts---5,654~samples
were retained (Table~\ref{tab:cohort}).

Healthy controls comprised 3,403~samples. Cancer cases comprised
2,251~samples: breast cancer ($n=1{,}280$), colorectal cancer (CRC;
$n=874$), prostate cancer ($n=93$; treated as exploratory throughout),
and lung cancer ($n=4$; retained only in binary pan-cancer
classification). Clinical stage annotations were recoverable for
731~cancer patients, all from GSE115513 (Stage~I: $n=231$; II:
$n=181$; III: $n=219$; IV: $n=100$). No ethical approval was required
as only publicly available, de-identified data were used.

\begin{table}[H]
\centering
\caption{Cohort composition across the four GEO source datasets.
HC~=~healthy control; CRC~=~colorectal cancer.
$^\dagger$Other: prostate ($n=93$), lung ($n=4$).}
\label{tab:cohort}
\begin{tabular}{lrrrrr}
\toprule
Dataset & $N$ & HC & Breast & CRC & Other$^\dagger$ \\
\midrule
GSE73002  & 4,113 & 2,740 & 1,280 & 874 & -- \\
GSE115513 & 1,513 &   649 &    -- &  -- & 864 \\
GSE126094 &    20 &    10 &    -- &  -- &  10 \\
GSE101684 &     8 &     4 &    -- &  -- &   4 \\
\midrule
Total     & 5,654 & 3,403 & 1,280 & 874 &  97 \\
\bottomrule
\end{tabular}
\end{table}

\subsection{Patient-Specific Tripartite ceRNA Graph Construction}

For each patient $p$, a directed tripartite graph
$G_p = (V, E)$ was constructed where the vertex set $V$ comprises
three disjoint node types---circRNAs, miRNAs, and mRNAs---and the
edge set $E$ encodes experimentally validated regulatory
relationships. The graph is patient-specific: active edges in $G_p$
depend on which nodes are co-expressed above the 25th-percentile
co-expression threshold computed within that patient's own expression
distribution.

miRNA--mRNA regulatory interactions were sourced from miRTarBase
v9.0~\citep{huang2022}, retaining only entries supported by strong
direct experimental evidence: CLIP-Seq, reporter assay, Western blot,
or quantitative PCR. This filtering yielded 10,770~high-confidence
edges. circRNA--miRNA sponge interactions were sourced from
CircInteractome~\citep{dudekula2016}. For datasets lacking direct mRNA
profiling (GSE115513, GSE126094, GSE101684), mRNA expression was
imputed conservatively from the miRNA--mRNA co-expression network
defined by miRTarBase: for each active miRNA, all validated mRNA
targets were included with a binary presence indicator.

Co-expression threshold sensitivity analysis confirmed AUROC variation
of no more than $\Delta=0.006$ across the 10th--40th percentile range,
establishing robustness to this design choice.

\subsection{Feature Extraction}

Twenty-one features were computed per patient (Table~\ref{tab:features}).
The four \emph{expression features} capture RNA abundance statistics:
\texttt{n\_circ\_expressed}, \texttt{mean\_expression},
\texttt{max\_expression}, and \texttt{expression\_std}, all computed on
$\log_2(\text{CPM}+1)$ transformed count values.

The seventeen \emph{topology features} capture structural properties of
$G_p$ across six categories: scale (\texttt{n\_nodes},
\texttt{n\_edges}), degree statistics (\texttt{mean\_degree},
\texttt{max\_degree}, \texttt{top5\_hub\_degree},
\texttt{top1\_circ\_degree}), centrality
(\texttt{max\_circ\_pagerank}~\citep{kolaczyk2009},
\texttt{top5\_hub\_betweenness}), community structure
(\texttt{community\_count} by Louvain algorithm,
\texttt{modularity} by Newman--Girvan criterion~\citep{newman2006}),
connectivity (\texttt{diameter}), information-theoretic measures
(\texttt{graph\_entropy}, \texttt{spectral\_entropy}), and TDA
features (\texttt{betti\_0}, \texttt{betti\_1},
\texttt{persistence\_entropy}~\citep{otter2017}).

\begin{table}[H]
\centering
\caption{Complete feature set. The Batch-Sensitive column indicates
whether the feature shows significant cross-dataset distributional
shift ($p<0.001$, Kruskal--Wallis). Cancer direction ($\uparrow$
higher in cancer, $\downarrow$ lower in cancer) is derived from
univariate effect-size analysis and SHAP importance.}
\label{tab:features}
\small
\begin{tabular}{llp{5.5cm}cc}
\toprule
Group & Feature & Description & Cancer & Batch-Sensitive \\
\midrule
\multirow{4}{*}{Expression}
 & mean\_expression    & Mean $\log_2(\text{CPM}+1)$           & $\downarrow$ & No  \\
 & max\_expression     & Max $\log_2(\text{CPM}+1)$            & $\downarrow$ & No  \\
 & expression\_std     & SD of $\log_2(\text{CPM}+1)$          & $\downarrow$ & No  \\
 & n\_circ\_expressed  & circRNAs above threshold              & $\downarrow$ & No  \\
\midrule
\multirow{2}{*}{Scale}
 & n\_nodes & Total nodes in $G_p$   & $\downarrow$ & \textbf{Yes} \\
 & n\_edges & Total edges in $G_p$   & $\downarrow$ & \textbf{Yes} \\
\midrule
\multirow{4}{*}{Degree}
 & mean\_degree       & Mean degree across all nodes          & $\downarrow$ & \textbf{Yes} \\
 & max\_degree        & Maximum degree                        & $\downarrow$ & \textbf{Yes} \\
 & top5\_hub\_degree  & Mean degree of top-5 hub nodes        & $\downarrow$ & \textbf{Yes} \\
 & top1\_circ\_degree & Degree of highest-degree circRNA      & $\downarrow$ & No  \\
\midrule
\multirow{2}{*}{Centrality}
 & max\_circ\_pagerank      & Max PageRank among circRNA nodes      & $\downarrow$ & No  \\
 & top5\_hub\_betweenness   & Mean betweenness of top-5 hubs        & $\downarrow$ & \textbf{Yes} \\
\midrule
Clustering & avg\_clustering & Mean clustering coefficient           & $\downarrow$ & No  \\
\midrule
\multirow{2}{*}{Community}
 & community\_count & Louvain community count               & $\uparrow$   & \textbf{Yes} \\
 & modularity       & Newman--Girvan modularity $Q$         & $\downarrow$ & \textbf{Yes} \\
\midrule
Connectivity & diameter & Longest shortest path in LCC          & $\uparrow$   & \textbf{Yes} \\
\midrule
\multirow{2}{*}{Information}
 & graph\_entropy    & Shannon entropy of degree distribution & varies & \textbf{Yes} \\
 & spectral\_entropy & Entropy of normalised Laplacian spectrum& varies & \textbf{Yes} \\
\midrule
\multirow{3}{*}{TDA}
 & betti\_0             & Connected components ($H_0$)          & $\uparrow$   & \textbf{Yes} \\
 & betti\_1             & Independent regulatory cycles ($H_1$) & varies & \textbf{Yes} \\
 & persistence\_entropy & Entropy of persistence diagram        & varies & No  \\
\bottomrule
\multicolumn{5}{l}{\small Batch-sensitive: 12/21 features (57.1\%). Non-sensitive: 9/21 features (42.9\%).}
\end{tabular}
\end{table}

\subsection{Degree-Preserving Null Network Validation}

To distinguish genuine higher-order structural signal from properties
of the degree sequence, a degree-preserving random edge permutation was
performed on a stratified subsample of 500~patients (250~cancer,
250~healthy) drawn from fold~1. Edges were randomly rewired 100~times
using the Maslov--Sneppen algorithm~\citep{maslov2002}, which preserves
exact in- and out-degrees while destroying community membership, path
lengths, cycle structure, and clustering. Topology-only XGBoost AUROC
on real graphs: 0.655; on degree-preserving permuted graphs: 0.531
($p < 0.001$, 1,000-iteration permutation test). This 12.4-percentage-
point reduction confirms that topology features encode structural
information beyond the degree sequence within the tested subsample.
Extension of this validation to all folds is pre-registered as follow-up
work.

\subsection{Preprocessing and Batch Harmonisation}

Raw counts were log-transformed as $x' = \log_2(\text{CPM}+1)$. All
21~features were $z$-score normalised using parameters estimated
exclusively from the training partition within each cross-validation
fold, preventing leakage of test-data statistics into feature scaling.
ComBat~\citep{johnson2007} was applied to the four expression features
only, without including the cancer/healthy label as a biological
covariate. A control experiment including the cancer label in ComBat
inflated AUROC to $>0.999$---a clear label-leakage signature---
confirming that supervised batch correction constitutes a form of
data leakage. In the LODO evaluation, ComBat was fitted exclusively on
the three training datasets and applied prospectively to the held-out
fourth dataset.

\subsection{Cross-Validation Protocols}

Two cross-validation protocols were employed to quantify the
performance gap induced by dataset-naive evaluation:

\paragraph{Protocol~1: Dataset-Naive CV (standard practice).}
Outer folds were constructed using \texttt{StratifiedKFold} ($k=5$),
stratified by class label only. Samples from the same GEO source
therefore appear in both training and test partitions within the outer
loop. Inner 3-fold cross-validation with 50~Optuna~TPE
trials~\citep{akiba2019} optimised XGBoost~\citep{chen2016}
hyperparameters. This protocol mirrors standard practice in the ceRNA
classifier literature.

\paragraph{Protocol~2: Dataset-Stratified CV (recommended practice).}
Outer folds were constructed using \texttt{StratifiedGroupKFold}
($k=4$, one per dataset), using the GEO dataset accession as the
grouping variable. This ensures entire datasets are held out in each
outer fold, preventing the model from exploiting platform-specific
batch signals. Split-local $z$-score normalisation (computed only on
training statistics) and platform covariates (sequencing technology
dummy variables: microarray vs.\ RNA-seq) were incorporated to isolate
biological from technical signal. The same inner Optuna protocol was
applied within each outer training fold.

Out-of-fold predicted probabilities from all outer folds were
concatenated before AUROC computation---the statistically correct
estimator under nested cross-validation. All 95\%~confidence intervals
were computed using 2,000~stratified bootstrap resamples with the
percentile method. Statistical significance was assessed using the
DeLong non-parametric test~\citep{delong1988} ($\alpha=0.05$). SHAP
TreeExplainer~\citep{lundberg2017} was used to compute global feature
importance on the full-dataset model trained under Protocol~1. LODO
validation followed the same ComBat and Optuna protocol, with all
harmonisation and tuning applied exclusively within training data.

All analyses were implemented in Python~3.10 (scikit-learn~1.3,
XGBoost~1.7, Optuna~3.2, SHAP~0.42, NetworkX~3.1, Ripser~0.6,
neuroCombat~0.9).

\subsection{Topology Sensitivity Analysis}

To quantify which graph-theoretic features encode dataset-specific
batch effects versus biological signal, a Kruskal--Wallis
non-parametric test~\citep{kruskal1952} was applied to each of the
21~features across the four GEO datasets. Features with $p < 0.001$
were classified as batch-sensitive. The reference dataset for
directional comparison was GSE73002, which provides genuine matched
mRNA profiles via exoRBase~\citep{li2018}; the three remaining datasets
rely on miRTarBase-imputed mRNA indicators. This analysis provides a
reproducible diagnostic for any multi-cohort ceRNA study.

\subsection{Full-Matrix Expression Baseline: Feasibility Analysis}

A full-matrix elastic net classifier trained on the complete circRNA
expression space was investigated as an alternative baseline. Direct
comparison with the hybrid model was precluded by three structural
incompatibilities: (i)~alternative large-scale circRNA cohorts
(GSE164174, GSE113740) entirely lack the matched miRNA expression
profiles required for ceRNA graph construction; (ii)~GSE164174
represents a single-cancer binary task (gastric cancer
vs.\ healthy)---a fundamentally simpler problem than the pan-cancer
formulation; and (iii)~raw circRNA count matrices in these cohorts
exhibit sparsity exceeding 40\%, violating the quality-control
thresholds established for the tripartite pipeline. This data
bifurcation---large circRNA-only cohorts versus smaller
circRNA+miRNA cohorts---constitutes a critical infrastructure
bottleneck in the field and motivates the development of matched
multi-omic liquid biopsy resources.

% -------------------------------------------------------
% 3. RESULTS
% -------------------------------------------------------
\section{Results}

\subsection{Feature-Level Discriminability and Batch Sensitivity}

Among expression features, \texttt{mean\_expression},
\texttt{max\_expression}, and \texttt{expression\_std} all showed
Cohen's $d > 1.0$ between cancer and healthy samples, confirming the
well-established strong differential expression signal in circRNA
liquid biopsy data. Expression features showed no significant
cross-dataset distributional shift ($p > 0.05$, Kruskal--Wallis),
confirming that ComBat harmonisation successfully corrected
expression-level batch effects.

Among topology features, 12 of 21 (57.1\%) showed significant
cross-dataset distributional shift ($p < 0.001$), including
\texttt{community\_count}, \texttt{modularity}, \texttt{diameter},
\texttt{spectral\_entropy}, \texttt{betti\_0}, \texttt{betti\_1},
\texttt{n\_nodes}, \texttt{n\_edges}, \texttt{mean\_degree},
\texttt{max\_degree}, \texttt{top5\_hub\_betweenness}, and
\texttt{graph\_entropy} (Table~\ref{tab:features};
Figure~\ref{fig:sensitivity}). The features with the largest
univariate cancer effect sizes---\texttt{community\_count}
(Cohen's $d=0.61$), \texttt{modularity} ($d=-0.55$), and
\texttt{diameter} ($d=0.48$)---are simultaneously among the most
batch-sensitive, indicating that their apparent discriminative power
conflates genuine regulatory fragmentation with platform-specific
graph construction artefacts.

\subsection{Performance Gap: Dataset-Naive vs.\ Dataset-Stratified CV}

Table~\ref{tab:main_results} presents the central finding of this
study. Under Protocol~1 (dataset-naive CV), the combined
topology--expression model achieved AUROC~0.927 (95\%~CI:
0.921--0.934) with sensitivity 66.3\% at 95\%~specificity.
Under Protocol~2 (dataset-stratified CV), AUROC ranged from 0.45 to
0.76 depending on the held-out dataset, a reduction of up to 0.48
AUROC points. The magnitude of this drop directly reflects the extent
to which Protocol~1 allows the model to exploit platform-specific
batch signals in the training data.

The expression-only baseline (4~summary statistics) achieved
AUROC~0.909 under Protocol~1. The hybrid model's apparent advantage
of $\Delta$AUROC~$=+0.018$ over this baseline (DeLong $p < 0.05$)
should be interpreted with caution given the inflated nature of both
estimates under Protocol~1.

Under Protocol~2, the dataset-stratified AUROC of 0.45 when
GSE115513 is held out and 0.76 when other datasets are held out
reflects the genuine cross-platform generalisation difficulty. The
model trained on three platforms and tested on a fourth encounters
batch effects that ComBat cannot fully resolve because differences in
RNA extraction, library preparation, and sequencing depth alter the
effective co-expression landscape and resulting graph topology.

\begin{table}[H]
\centering
\caption{Central result: performance comparison under dataset-naive
(Protocol~1) versus dataset-stratified (Protocol~2) cross-validation.
Sensitivity values at 95\%~specificity. LODO = leave-one-dataset-out.
$^\ddagger$Protocol~2 AUROC varies by held-out dataset; range reported.}
\label{tab:main_results}
\begin{tabular}{llcc}
\toprule
Protocol & Model & AUROC (95\% CI) & Sens.\@ 95\% Sp. \\
\midrule
\multirow{3}{*}{\shortstack[l]{Protocol~1\\(Dataset-naive CV)}}
 & Expression-only (4 features) & 0.909 (0.901--0.916) & 55.6\% \\
 & Topology-only (17 features)  & 0.655 (0.640--0.670) & 24.3\% \\
 & Hybrid (21 features)         & \textbf{0.927 (0.921--0.934)} & \textbf{66.3\%} \\
\midrule
\multirow{1}{*}{\shortstack[l]{Protocol~2\\(Dataset-stratified CV)}}
 & Hybrid (21 features) & 0.45--0.76$^\ddagger$ & -- \\
\midrule
\multirow{2}{*}{LODO}
 & GSE73002 held out ($n=4{,}113$)  & 0.749 (0.733--0.762) & -- \\
 & GSE115513 held out ($n=1{,}513$) & 0.540 (0.510--0.571) & -- \\
\bottomrule
\multicolumn{4}{l}{\small Topology-only (perm.): AUROC~0.531 ($p < 0.001$ vs.\ real; 500-patient subsample).}
\end{tabular}
\end{table}

\subsection{LODO External Validation}

LODO results (Table~\ref{tab:main_results}) confirm the cross-platform
generalisation difficulty identified by Protocol~2. When GSE73002
($n=4{,}113$) was withheld, the model achieved AUROC~0.749
(95\%~CI: 0.733--0.762), demonstrating that a cancer detection signal
does transfer across platforms when the training set is large and
diverse. When GSE115513 ($n=1{,}513$) was withheld---the only dataset
with stage annotations, generated under a single-institution uniform
protocol---the model achieved near-chance AUROC~0.540
(95\%~CI: 0.510--0.571) despite unsupervised ComBat harmonisation.
GSE126094 ($n=20$) and GSE101684 ($n=8$) produced statistically
uninformative confidence intervals due to insufficient sample size and
are reported for completeness only.

The LODO failure on GSE115513 is not random. Unsupervised ComBat
corrects additive expression shifts but cannot compensate for
differences in RNA extraction and circRNA enrichment protocols that
alter the effective co-expression landscape and resulting graph
topology---precisely the batch-sensitive features identified in
Section~3.1. This finding is consistent with known challenges in
multi-cohort circRNA meta-analyses and is accurately predicted by the
topology sensitivity analysis.

\subsection{Stage Stratification: Within-Protocol Evidence Only}

Within the GSE115513 staged sub-cohort ($n=731$~cancer,
$n=649$~healthy; single institution, uniform protocol), Stage~I AUROC
was 0.924 for the hybrid model versus 0.886 for the expression-only
baseline, translating to a 2.46-fold improvement in Stage~I
sensitivity at 95\%~specificity: 53.2\% vs.\ 21.6\%. The
Kruskal--Wallis test across Stages~I--IV yielded $p=0.31$, confirming
that model performance does not depend on late-stage disease burden.

\textbf{Scope caveat}: These results are within-cohort estimates
derived from a single uniformly processed dataset and must not be
interpreted as cross-platform clinical performance. The LODO failure
of GSE115513 (AUROC~0.540) means the Stage~I sensitivity improvement
has not been validated externally and likely reflects single-platform
internal consistency rather than generalizable biological signal.

\subsection{SHAP Feature Importance and Biological Interpretability}

Global SHAP feature importance (Figure~\ref{fig:shap}) computed under
Protocol~1 shows expression features dominating by mean absolute SHAP
value, consistent with the strong expression-only baseline. Among
topology features, \texttt{community\_count}
(mean~$|\text{SHAP}|=0.305$) is most important, followed by
\texttt{diameter} (0.249), \texttt{mean\_degree} (0.152), and
\texttt{modularity} (0.125).

\textbf{Critical interpretation}: The topology features with the
highest SHAP importance---\texttt{community\_count},
\texttt{diameter}, \texttt{mean\_degree}, \texttt{modularity}---are
all batch-sensitive ($p < 0.001$; Table~\ref{tab:features}). Their
SHAP importance under Protocol~1 therefore cannot be cleanly
attributed to cancer biology; it reflects a mixture of biological
regulatory fragmentation signal and platform-specific graph
construction differences. This observation does not invalidate the
null network validation result (topology-only AUROC drops 12.4~pp
under degree-preserving permutation), which confirms that genuine
structural signal exists in the data. It does, however, mean that the
SHAP importance rankings should not be used to rank circRNA biomarkers
for clinical follow-up without first replicating them in a
platform-controlled study.

% -------------------------------------------------------
% 4. DISCUSSION
% -------------------------------------------------------
\section{Discussion}

\subsection{The Dataset-Naive CV Problem in ceRNA Liquid Biopsy Research}

The central finding of this study is that dataset-naive cross-validation
inflates apparent AUROC by up to 0.48~points in a multi-cohort ceRNA
liquid biopsy meta-cohort. This gap is not a statistical curiosity; it
has direct consequences for how the field interprets published results.
A classifier reporting AUROC~0.93 under Protocol~1 may be doing little
more than learning to distinguish Illumina miRNA-seq from exosome
circRNA microarray---a platform classification task with no clinical
utility.

The mechanism is straightforward. ceRNA graph topology features are
constructed by thresholding expression values, meaning the set of
active nodes and edges depends directly on the underlying expression
distribution. When that distribution carries platform-specific
signatures---as 57.1\% of topology features demonstrably do---the
topology features inherit those signatures. A model trained on three
platforms that sees samples from all four in cross-validation can
learn these signatures implicitly, even without explicit platform
labels.

This problem is not unique to ceRNA classifiers. It is a specific
instance of the broader dataset shift problem in multi-source
biomedical machine learning~\citep{leek2010}. What is specific to this
subfield is that the problem is compounded by the network construction
step: batch effects enter not only through raw expression values but
also through the graph topology derived from those values, creating a
two-stage amplification of platform signal.

\subsection{What the Honest Numbers Mean}

The dataset-stratified AUROC of 0.45--0.76 and LODO AUROC of 0.540--0.749
represent the genuine cross-platform performance of the framework.
The wide range reflects the heterogeneity of the four GEO datasets in
terms of cancer type, sequencing platform, and sample size. The upper
bound (0.749, GSE73002 LODO) demonstrates that a genuine pan-cancer
ceRNA signal does exist and transfers across platforms when training data
is large and diverse. The lower bound (0.540, GSE115513 LODO) demonstrates
that this signal is not robust to all platform combinations, and that
the dataset-stratified CV range of 0.45--0.76 accurately anticipates
this heterogeneity.

The degree-preserving null model result (topology-only AUROC drops from
0.655 to 0.531, $p<0.001$) confirms that genuine structural signal
exists independent of the degree sequence---ceRNA network topology does
encode cancer-associated regulatory fragmentation. The problem is not
that topology is uninformative; it is that topology features are
confounded with platform identity to a degree that makes their
discriminative contribution difficult to isolate without a controlled
experimental design.

\subsection{Implications for the Field}

Prior multi-cohort ceRNA classifier studies that employed dataset-naive
cross-validation likely overestimate their performance. This does not
mean those studies are invalid; it means their reported AUROC values
should be interpreted as upper bounds on performance under idealised
(within-platform) conditions rather than as estimates of
cross-platform clinical utility.

Three practical recommendations follow from this work. First, any
multi-cohort study should report both dataset-naive and
dataset-stratified CV results, and should treat the latter as the
primary performance estimate. Second, SHAP importance rankings derived
from dataset-naive models should not be used to rank biomarkers for
clinical follow-up when the top-ranked features are batch-sensitive.
Third, the topology sensitivity diagnostic introduced here---
Kruskal--Wallis cross-dataset distributional testing of all graph
features---should be run as a standard pre-analysis step in any
multi-cohort ceRNA network study.

\subsection{Limitations}

\begin{enumerate}
  \item Dataset-stratified AUROC range (0.45--0.76) is based on four
  datasets; the range may narrow or widen with additional cohorts.

  \item Degree-preserving null validation was executed on 500~patients
  from fold~1 only (8.8\% of the full cohort) due to computational
  cost; extension to all folds is pre-registered as follow-up analysis.

  \item Imputed mRNA data (miRTarBase-derived binary indicators) for
  three of four datasets introduces structural assumptions that differ
  fundamentally from the matched mRNA profiles available via exoRBase
  for GSE73002. This imputation strategy likely contributes to the
  batch sensitivity observed in topology features.

  \item Stage-specific results are within-cohort estimates on
  GSE115513 only; cross-platform Stage~I performance remains to be
  established.

  \item The prostate cancer sub-cohort ($n=93$) is insufficient for
  robust inference; all prostate results are exploratory.

  \item A full-matrix elastic net expression baseline evaluated on
  matched data was not feasible due to the structural bifurcation
  between circRNA-only and circRNA+miRNA public cohorts, as described
  in Section~2.7.
\end{enumerate}

% -------------------------------------------------------
% 5. CONCLUSION
% -------------------------------------------------------
\section{Conclusion}

Dataset-naive cross-validation inflates apparent AUROC by up to
0.48~points in multi-cohort ceRNA liquid biopsy classifiers, because
57.1\% of ceRNA graph topology features encode platform-specific batch
effects alongside biological signal. Genuine cross-platform performance,
measured by dataset-stratified CV and LODO validation, falls in the
range AUROC~0.45--0.76---substantially below the 0.90--0.95 values
routinely reported in the literature. A degree-preserving null network
validation confirms that cancer-associated structural signal genuinely
exists in ceRNA topology, but this signal is not currently separable
from platform identity without standardised pre-analytical protocols.
The topology sensitivity diagnostic and dataset-stratified evaluation
pipeline introduced here provide reproducible tools for the field to
benchmark future classifiers honestly, and motivate prospective
multi-site studies with harmonised sequencing protocols as the path
toward clinically credible ceRNA network-based cancer detection.

% -------------------------------------------------------
% DATA AVAILABILITY
% -------------------------------------------------------
\section*{Data Availability}

All datasets are publicly available from NCBI GEO
(\url{https://www.ncbi.nlm.nih.gov/geo/}) under accession numbers
GSE73002, GSE115513, GSE126094, and GSE101684. Processed feature
matrices, the topology sensitivity diagnostic script, and the
dataset-stratified cross-validation pipeline are available at
\url{https://github.com/Sharon-codes/Cancer-Research-JBI}.

% -------------------------------------------------------
% REFERENCES
% -------------------------------------------------------
\bibliography{reference}

% -------------------------------------------------------
% FIGURES  (legends here; image files uploaded separately)
% -------------------------------------------------------
\newpage
\section*{Figure Legends}

\begin{description}

\item[Figure 1.] \textbf{Topology sensitivity analysis across four GEO
datasets.} Box plots showing the cross-dataset distribution of
representative topology features: (A)~\texttt{community\_count},
(B)~\texttt{modularity}, (C)~\texttt{diameter},
(D)~\texttt{spectral\_entropy}, (E)~\texttt{betti\_0}, and
(F)~\texttt{mean\_expression} (non-batch-sensitive control). Kruskal--Wallis
$p$-values are annotated above each panel. Features in panels A--E show
significant cross-dataset distributional shift ($p < 0.001$), indicating
batch sensitivity. The expression control in panel~F shows no significant
shift ($p > 0.05$), confirming that ComBat harmonisation successfully
corrected expression-level batch effects while topology features
remain confounded.

\item[Figure 2.] \textbf{SHAP beeswarm plot for the hybrid model
(Protocol~1, dataset-naive CV).} Features are ranked by mean absolute
SHAP value. Colour encodes feature value (red~=~high, blue~=~low).
Features marked with~$\star$ in the y-axis label are batch-sensitive
($p < 0.001$, Kruskal--Wallis). The dominance of batch-sensitive topology
features in the top-10 SHAP rankings illustrates why Protocol~1
importance scores cannot be used to rank circRNA biomarkers for clinical
follow-up without platform-controlled replication.

\item[Figure 3.] \textbf{Structural comparison of representative
patient-specific ceRNA subgraphs.} Left: Healthy patient subgraph
exhibiting dense modular organisation with high-degree circRNA hubs.
Right: Stage~I cancer patient subgraph exhibiting hub dissolution,
reduced modularity, and network fragmentation. Nodes are coloured by
type: green~=~circRNA, orange~=~miRNA, grey~=~mRNA.

\item[Figure 4.] \textbf{Performance comparison across evaluation
protocols.} Bar chart comparing AUROC under Protocol~1 (dataset-naive
CV), Protocol~2 (dataset-stratified CV, range shown as error bar), and
LODO validation (per held-out dataset). The gap between Protocol~1 and
Protocol~2/LODO quantifies the inflation attributable to dataset-naive
evaluation.

\end{description}

\end{document}